\section*{Bab \ref{ch:b3:climate-models} --- Model Iklim}

\begin{solution}{exo:b3:climate-models:1}
$A = 14.0 - 1.5 = \num{12.5}$ dan $B = 0.14 - 0.05 = \num{0.09}$, dalam
\unit{kcal.cm^{-2}} per bulan. Satu bulan adalah \qty{2.63e6}{s}, sehingga
\qty{1}{kcal.cm^{-2}} per bulan adalah $\num{4.184e7}/\num{2.63e6} =
\qty{15.91}{W/m^2}$: $A = \qty{199}{W/m^2}$, $B =
\qty{1.43}{W.m^{-2}.\celsius^{-1}}$. Pada \qty{15}{\celsius}, $199 + 21.5 =
\qty{220}{W/m^2}$.
\end{solution}

\begin{solution}{exo:b3:climate-models:2}
$(238.0 - 211.9)/1.816 = \qty{14.4}{\celsius}$; dengan 0.31, $(234.6 -
211.9)/1.816 = \qty{12.5}{\celsius}$. Setiap 0.01 albedo setara dengan
\qty{3.4}{W/m^2} dan $3.4/1.816 = \qty{1.9}{\celsius}$.
\end{solution}

\begin{solution}{exo:b3:climate-models:3}
$C = 1025\times3992\times60 = \qty{2.46e8}{J.m^{-2}.K^{-1}}$; dibagi
dengan \qty{3.156e7}{s} dalam setahun, \qty{7.78}{W.yr.m^{-2}.K^{-1}}. Lalu $\tau
= 7.78/1.816 = \num{4.3}\text{ tahun}$, dan $7.78/1.33 = \num{5.8}\text{ tahun}$.
\end{solution}

\begin{solution}{exo:b3:climate-models:4}
ECS $= 3.93/1.33 = \qty{2.95}{\celsius}$; TCR $\approx 3.93/1.95 =
\qty{2.02}{\celsius}$; rasio \num{0.68}.
\end{solution}

\begin{solution}{exo:b3:climate-models:5}
$Q = \qty{335}{W/m^2}$. Bola salju: $(335\times0.38 - 211.9)/1.816 =
\qty{-46.6}{\celsius}$; hangat: $(335\times0.70 - 211.9)/1.816 =
\qty{12.4}{\celsius}$, masih di atas \qty{10}{\celsius}; tanjakan: $335(0.54 +
0.016\,T) = 211.9 + 1.816\,T$ memberikan $T = \qty{8.7}{\celsius}$. $N' =
\qty{-1.816}{W.m^{-2}.\celsius^{-1}}$ pada dua yang pertama, stabil, dan $-1.816
+ 335\times0.016 = \num{+3.54}$ pada tanjakan, tidak stabil.
\end{solution}

\begin{solution}{exo:b3:climate-models:6}
Cabang hangat berakhir ketika suhunya mencapai \qty{0}{\celsius}:
$0.70\,Q = A$, $Q = \qty{302.7}{W/m^2}$, konstanta surya sebesar
\qty{1211}{W/m^2}, \qty{11.0}{\percent} di bawah nilai hari ini. Bola salju mencair
saat $0.38\,Q = A - 20B$, $Q = \qty{462.1}{W/m^2}$, $S_0 =
\qty{1848}{W/m^2}$. Ambang batasnya adalah $B(T^* - T_F)/(Q_0(1 - \alpha_F))$,
yaitu selisih antara suhu hari ini $T^*$ dan suhu saat es
mulai menyebar, dikonversi menjadi sinar Matahari: $1.816\times14.37/238 =
\qty{11.0}{\percent}$ di sini, $1.816\times4.37/238 = \qty{3.3}{\percent}$ dalam
teks.
\end{solution}

\begin{solution}{exo:b3:climate-models:7}
(a)~\qty{-45.5}{\celsius}. (b)~$340\times0.38 + F = 211.9 - 18.2$: $F =
\qty{64.5}{W/m^2}$, 16 kali gaya radiatif akibat penggandaan. (c)~Selama
gaya radiatif berlangsung, planet melompat ke cabang hangat, $(238.0 + 64.5 -
211.9)/1.816 = \qty{50}{\celsius}$; setelah gaya itu dihapus, planet mendingin
sepanjang cabang hangat hingga \qty{14.4}{\celsius}, yang ada pada Matahari hari ini:
ia tidak kembali ke bola salju.
\end{solution}

\begin{solution}{exo:b3:climate-models:8}
$\Sigma = 1.80/7.0 + 0.70/90 = \qty{0.2649}{yr^{-1}}$, $\Pi = 1.10\times0.70/
(7.0\times90) = \qty{1.222e-3}{yr^{-2}}$. Akar-akarnya memberikan $\tau_f =
\num{3.84}\text{ tahun}$ dan $\tau_s = \num{213}\text{ tahun}$; pendekatannya,
$7.0/1.80 = \num{3.89}\text{ tahun}$, $90\times1.80/(1.10\times0.70) =
\num{210}\text{ tahun}$. $a_f = 0.597$ tepat, $1.10/1.80 = 0.611$
kira-kira. ECS $= 3.93/1.10 = \qty{3.57}{\celsius}$; TCR $\approx
3.93/1.80 = \qty{2.18}{\celsius}$.
\end{solution}

\begin{solution}{exo:b3:climate-models:9}
Penyebut $B + \beta = \num{5.555}$. Khatulistiwa: $(340\times1.221\times0.70 -
211.9 + 3.74\times12.81)/5.555 = \qty{22.8}{\celsius}$; kutub:
$(340\times0.506\times0.38 - 211.9 + 47.9)/5.555 = \qty{-17.7}{\celsius}$.
Transpor menghilangkan $3.74\times(22.8 - 12.8) = \qty{37}{W/m^2}$ di
khatulistiwa dan membawa $3.74\times30.6 = \qty{114}{W/m^2}$ ke kutub.
\end{solution}

\begin{solution}{exo:b3:climate-models:10}
(a)~$(1 - x^2)P_2' = 3x - 3x^3$, yang turunannya adalah $3 - 9x^2 = -6P_2$.
(b)~Suku transpornya adalah $-\mathcal D T_2(-6P_2) = 6\mathcal DT_2P_2$.
Keseimbangan $Q(1 + s_2P_2)(1 - \alpha) - A - B(T_0 + T_2P_2) - 6\mathcal
DT_2P_2 = 0$ berlaku untuk semua $x$ jika bagian konstanta dan bagian $P_2$-nya
lenyap, yang memberikan $T_0$ dan $T_2$. (c)~$T(0) - T(1) = T_2(-\tfrac12 - 1)
= -1.5\,T_2$. Tanpa transpor $T_2 = 340\times(-0.476)\times0.70/1.816 =
\qty{-62.4}{\celsius}$: selisih sebesar \qty{94}{\celsius}. (d)~$T_2 =
-55.3/1.5 = \qty{-36.9}{\celsius}$, sehingga $B + 6\mathcal D = 113.3/36.9 =
3.07$ dan $\mathcal D = \qty{0.21}{W.m^{-2}.\celsius^{-1}}$.
\end{solution}

\begin{solution}{exo:b3:climate-models:11}
(a)~$r = \tfrac12\Sigma\big(1 \pm \sqrt{1 - 4\Pi/\Sigma^2}\big)$; dengan
$\sqrt{1 - \epsilon} \approx 1 - \epsilon/2$, akarnya adalah $\Sigma -
\Pi/\Sigma \approx \Sigma$ dan $\Pi/\Sigma$. (b)~Dengan $T_d = 0$, $C\,\mathrm
dT/\mathrm dt = kt - (\gamma - \lambda)T$, yang solusinya dari keadaan diam adalah $T =
k(t - \tau_f)/(\gamma - \lambda) + k\tau_f\mathrm e^{-t/\tau_f}/(\gamma -
\lambda)$, $\tau_f = C/(\gamma - \lambda)$. Untuk $t \gg \tau_f$, $T \approx
kt/(\gamma - \lambda)$, dan pada saat penggandaan, $kt = F_{2\times}$: TCR
$\approx F_{2\times}/(\gamma - \lambda)$. (c)~$\Sigma = 1.95/8.1 +
0.62/110 = \qty{0.2464}{yr^{-1}}$, $\Pi = 1.33\times0.62/(8.1\times110) =
\qty{9.25e-4}{yr^{-2}}$; $\tau_f = \num{4.12}\text{ tahun}$ dan $\tau_s =
\num{262}\text{ tahun}$ (kira-kira $1/\Sigma = \num{4.06}\text{ tahun}$ dan
$\Sigma/\Pi = \num{266}\text{ tahun}$); $4\Pi/\Sigma^2 = 0.061$.
\end{solution}

\begin{solution}{exo:b3:climate-models:12}
(a)~Rata-rata \qty{3.50}{\celsius}, simpangan baku \qty{0.81}{\celsius}.
(b)~ECS $= 1.88 + 1.86\,X$; korelasi \num{0.992}; simpangan baku
residu \qty{0.11}{\celsius}. (c)~$1.88 + 1.86\times0.7 =
\qty{3.18}{\celsius}$, dengan sebaran $\sqrt{(1.86\times0.1)^2 + 0.11^2} =
\qty{0.22}{\celsius}$, kira-kira seperempat dari sebaran ansambel.
(d)~Tidak ada mekanisme fisis yang menghubungkan $X$ dengan umpan balik; $X$ yang
ditemukan dengan mencoba banyak observabel, sehingga korelasinya mungkin hanya kebetulan; bias
yang dimiliki bersama oleh kedelapan model; model dari keluarga yang sama dihitung sebagai
independen.
\end{solution}

\begin{solution}{pb:b3:climate-models:1}
\textbf{1.} $Q_0 = 1360/4 = \qty{340}{W/m^2}$: planet menangkap
sinar Matahari pada cakram berluas $\pi R^2$ dan menyebarkannya ke bola berluas
$4\pi R^2$.

\textbf{2.} $340\times0.70 = \qty{238.0}{W/m^2}$; $211.9 + 1.816\times15 =
\qty{239.1}{W/m^2}$.

\textbf{3.} $(238.0 - 211.9)/1.816 = \qty{14.4}{\celsius}$.

\textbf{4.} $(129.2 - 211.9)/1.816 = \qty{-45.5}{\celsius}$, di bawah
\qty{-10}{\celsius}.

\textbf{5.} $340(0.54 + 0.016\,T) = 211.9 + 1.816\,T$: $T =
(183.6 - 211.9)/(1.816 - 5.44) = \qty{7.8}{\celsius}$.

\textbf{6.} $(238.0/\num{5.670e-8})^{1/4} = \qty{254.5}{K}$,
\qty{-18.6}{\celsius}: \qty{33}{\celsius} di bawah kesetimbangan hangat, yakni
efek rumah kaca planet yang termuat dalam $A$ dan $B$.

\textbf{7.} $N' = -B = \qty{-1.816}{W.m^{-2}.\celsius^{-1}}$ pada
bola salju dan kesetimbangan hangat, keduanya stabil; $-1.816 + 5.44 =
\num{+3.62}$ pada yang di tengah, tidak stabil.

\textbf{8.} $\tau = 8.1/3.62 = \num{2.2}\text{ tahun}$.

\textbf{9.} $8.1/1.816 = \num{4.5}\text{ tahun}$.

\textbf{10.} $3.93/1.816 = \qty{2.2}{\celsius}$, di bawah taksiran terbaik yang dinilai,
yakni \qty{3}{\celsius} dan di bawah rentangnya yang \emph{mungkin}, 2.5 sampai
\qty{4}{\celsius}: $B$ pada model, yang dicocokkan di seluruh lintang, bukan
parameter umpan balik iklim secara keseluruhan.

\textbf{11.} $\mathrm e^{-t/\tau} = 0.05$ pada $t = 3\tau =
\num{13}\text{ tahun}$. Iklim sebenarnya juga memanaskan samudra dalam, yang
skala waktunya berabad-abad.

\textbf{12.} Pada $T_F = \qty{10}{\celsius}$; $Q = (211.9 + 18.2)/0.70 =
\qty{328.7}{W/m^2}$.

\textbf{13.} $S_0 = \qty{1315}{W/m^2}$, \qty{3.3}{\percent} di bawah nilai hari ini.

\textbf{14.} Ke bola salju, $(328.7\times0.38 - 211.9)/1.816 =
\qty{-47.9}{\celsius}$.

\textbf{15.} $Q = (211.9 - 18.2)/0.38 = \qty{509.8}{W/m^2}$, $S_0 =
\qty{2039}{W/m^2}$, \qty{50}{\percent} di atas nilai hari ini.

\textbf{16.} Turun sepanjang cabang hangat hingga \qty{10}{\celsius} pada
\qty{1315}{W/m^2}, suhu turun ke \qty{-48}{\celsius}, lalu naik kembali sepanjang
cabang bola salju hingga \qty{-45.5}{\celsius} pada \qty{1360}{W/m^2}: planet itu
tetap beku. Hanya di atas \qty{2039}{W/m^2} ia akan melompat kembali ke
cabang hangat.

\textbf{17.} $\bar\alpha = 0.30\times0.9738 + 0.62\times0.0262 =
\num{0.3084}$; $\bar T = (340\times0.6916 - 211.9)/1.816 =
\qty{12.80}{\celsius}$.

\textbf{18.} $T_c = (340\times0.5643\times0.50 - 211.9 + 3.739\times12.80)/
5.555 = \qty{-12.26}{\celsius}$, \qty{3.5}{\celsius} lebih hangat daripada
reanalisis pada \ang{72}.

\textbf{19.} $5.555\times(-12.26) + 211.9\times3.059 = \num{580.1}$, dengan
$\beta/B = \num{2.059}$.

\textbf{20.} Penyebutnya 1.7053, 1.7473, 1.7565, 1.7488, dan
1.7362, serta $Q/Q_0 = \num{1.0005}$, \num{0.9764}, \num{0.9713},
\num{0.9756} dan \num{0.9827}.

\textbf{21.} Dekat \ang{47.5}, pada \num{0.971}. Garis es di sebelah kutub
dari minimum stabil, $\mathrm dQ/\mathrm dx_s > 0$; yang ke arah khatulistiwa
tidak stabil.

\textbf{22.} $Q/Q_0 = 1$ juga antara \ang{30} dan \ang{40}, dekat
\ang{30} (\ang{30.2} menurut interpolasi): garis es yang tidak stabil. Pada \ang{0},
$Q/Q_0 = 580.1/(0.6106 + 0.7824)/340 = \num{1.225}$: bola salju bertahan
untuk Matahari mana pun hingga \qty{22}{\percent} lebih terang daripada nilai hari ini, sehingga
keadaan itu mungkin terjadi hari ini.

\textbf{23.} \qty{2.9}{\percent}, dibanding \qty{3.3}{\percent} untuk
model nol-dimensi dan \qty{1.6}{\percent} untuk Budyko; es berada
dekat \ang{48} pada ambang, dibanding \ang{50} miliknya. Model lintang
membeku lebih cepat daripada yang global, karena sinar Matahari paling lemah di
tempat es terbentuk.

\textbf{24.} Planet model membeku bila penurunannya
\textbf{\qty[text-series-to-math=true]{2.9}{\percent}} dari konstanta surya
(\qty{2.87}{\percent} menurut perhitungan).
\end{solution}
